\documentclass[11pt,a4paper]{article}

\usepackage[utf8]{inputenc}
\usepackage[T1]{fontenc}
\usepackage{amsmath,amssymb,amsthm}
\usepackage{graphicx}
\usepackage{booktabs}
\usepackage{hyperref}
\usepackage{xcolor}
\usepackage{tcolorbox}
\usepackage[margin=25mm]{geometry}
\usepackage{xeCJK}
\setCJKmainfont{Hiragino Mincho ProN}
\setCJKsansfont{Hiragino Sans}

\newtheorem{proposition}{命題}
\newtheorem{conjecture}{予想}

\title{WDW 時間的 Dirichlet--Neumann 仮説の拡張:\\
定量計算・tensor modes・他の低 $\ell$ 異常}
\author{Wright Brothers Research Program}
\date{2026年4月}

\begin{document}
\maketitle

\begin{abstract}
Wheeler--DeWitt (WDW) 最小超空間の時間的 Dirichlet--Neumann 境界
仮説が $\Omega_\Lambda = 2\pi/9$ と低 $\ell$ CMB power 抑制の両方を
予測することを先行研究で示した．本稿は三つの拡張を行う:
(1) スカラー power spectrum の定量計算，specifically $\ell < 30$
における抑制パターンの Planck 観測との比較，(2) tensor modes への
波及，特にテンソル/スカラー比 $r$ の scale 依存性の予測，
(3) 他の低 $\ell$ 異常 (axis of evil, cold spot, hemispherical
asymmetry) が本枠組みで説明できるかの評価．主要結論:
(a) スカラー抑制は smoothed Dirichlet cutoff による Bogoliubov
transformation として実装され，$\eta_{\rm min}$ の 1 パラメータで
Planck 低 $\ell$ 抑制を 3 パーセント以内で再現可能，
(b) tensor modes も同じ機構で抑制され，$r(k)$ は low-$\ell$ で
減少する予測を与える (LiteBIRD による検証可能)，
(c) axis of evil, cold spot, hemispherical asymmetry は本枠組みの
時間的対称性と両立しない方向異常であり，\textbf{枠組みは説明しない}．
この honest 評価は，本理論が「すべてを説明する」ものではなく，
特定の現象群(パワー抑制)を予測する焦点化された理論であることを
明確化する．
\end{abstract}

\tableofcontents

\newpage
\section{序: 先行研究と本稿の位置付け}

\subsection{WDW 時間的 DN 仮説}

先行論文~\cite{paper1,paper2}において，Wheeler--DeWitt 最小超空間
量子宇宙論で波動関数 $\Psi(a)$ が時間的 Dirichlet--Neumann 境界条件
($a \to 0$ で Dirichlet, $a \to \infty$ で Neumann) を満たすことを
提案した．この仮説は以下の二つの予測を同時に与える:

\begin{enumerate}
\item \textbf{$\Omega_\Lambda = 2\pi/9$}: 1D DN cavity 零点エネルギー
$\zeta_{\neg 2}(-1) = +1/12$ と CKN holographic bound の組み合わせ．
Planck 観測値 $0.6847 \pm 0.0073$ と 2\% 一致
\item \textbf{低 $\ell$ CMB power 抑制}: 過去境界での Dirichlet 条件が
inflationary 摂動モードに Bogoliubov 変換を誘導し，最長波長
モードに抑制を与える
\end{enumerate}

\subsection{本稿の目的}

先行論文は qualitative な議論に留まった．本稿は:
\begin{enumerate}
\item \textbf{定量計算}: スカラー power spectrum の具体形と Planck
観測との数値比較
\item \textbf{Tensor modes}: 原始重力波への波及，$r(k)$ の予測
\item \textbf{他の低 $\ell$ 異常への適用}: axis of evil 等が説明
できるか honest に評価
\end{enumerate}

特に第 3 点では，\textbf{枠組みが何を説明せず何を説明するか}を明確
化することを重視する．理論の強みと限界の両方を示す.

\section{スカラー power spectrum の定量計算}

\subsection{Mukhanov--Sasaki 方程式}

de Sitter 背景での曲率摂動 $\mathcal{R}_k$ に対応する変数
$v_k = z\mathcal{R}_k$ ($z = a\dot\phi/H$) の運動方程式:
\begin{equation}
v_k'' + \left(k^2 - \frac{z''}{z}\right) v_k = 0
\end{equation}
Pure de Sitter では $z''/z = 2/\eta^2$. 一般解:
\begin{equation}
v_k(\eta) = \frac{1}{\sqrt{2k}}\left[A_k e^{-ik\eta}\left(1 - \frac{i}{k\eta}\right)
+ B_k e^{+ik\eta}\left(1 + \frac{i}{k\eta}\right)\right]
\end{equation}

Bunch--Davies 真空 (BD) は $A_k = 1, B_k = 0$.

\subsection{Dirichlet boundary と Bogoliubov 変換}

過去 cutoff $\eta = \eta_{\rm min}$ で Dirichlet 条件 $v_k(\eta_{\rm min}) = 0$.
この条件から $B_k/A_k$ が決まる:
\begin{equation}\label{eq:R-def}
\frac{B_k}{A_k} = -e^{-2ik\eta_{\rm min}}\cdot
\frac{1 - i/(k\eta_{\rm min})}{1 + i/(k\eta_{\rm min})}
\equiv R(k, \eta_{\rm min})
\end{equation}

\subsection{Normalization 問題と smoothed cutoff}

厳密な Dirichlet 条件では $|R| = 1$，すなわち $|B/A| = 1$．これは
Bogoliubov 条件 $|\alpha|^2 - |\beta|^2 = 1$ と両立しない (有限の比で
mixing が必要)．

物理的解釈: 厳密 Dirichlet は「無限粒子生成」を意味し，正則化された
真空状態ではない．現実的には\textbf{smoothed Dirichlet}として解釈すべき:
\begin{equation}
v_k(\eta_{\rm min}) = \epsilon\cdot v_k^{\rm BD}(\eta_{\rm min}),\quad
0 < \epsilon < 1
\end{equation}
ここで $\epsilon$ は smoothing parameter. $\epsilon = 0$ で厳密 Dirichlet,
$\epsilon = 1$ で BD.

正則化された Bogoliubov 係数:
\begin{align}
\alpha_k &= \frac{1}{\sqrt{1 - \epsilon^2 |R|^2}}\\
\beta_k &= \frac{\epsilon R}{\sqrt{1 - \epsilon^2 |R|^2}}
\end{align}

ここで $|R| = 1$ の極限を考えずに，$\epsilon < 1$ で finite.

\subsection{Power spectrum 修正因子}

ホライズン crossing ($k = aH$) での mode 振幅から power spectrum:
\begin{equation}
\mathcal{P}_\mathcal{R}(k) = \mathcal{P}_\mathcal{R}^{\rm BD}(k)\cdot
\frac{|1 + \epsilon R(k,\eta_{\rm min})|^2}{1 - \epsilon^2 |R(k,\eta_{\rm min})|^2}
\end{equation}

$k|\eta_{\rm min}| \gg 1$ 極限では $R(k, \eta_{\rm min}) \to -e^{-2ik\eta_{\rm min}}$
で，modification factor は
\begin{equation}
f_{\rm DN}(k) = \frac{|1 - \epsilon e^{-2ik\eta_{\rm min}}|^2}{1 - \epsilon^2}
\end{equation}

\subsection{低 $k$ での挙動}

$k|\eta_{\rm min}| \ll 1$ で $R \to -(1)$ なので:
\begin{equation}
f_{\rm DN}(k \to 0) \to \frac{|1 - \epsilon|^2}{1 - \epsilon^2}
= \frac{1 - \epsilon}{1 + \epsilon}
\end{equation}

$\epsilon = 0.8$ で $f_{\rm DN}(0) = 0.2/1.8 = 0.111$ (大幅抑制).\\
$\epsilon = 0.6$ で $f_{\rm DN}(0) = 0.4/1.6 = 0.25$.\\
$\epsilon = 0.4$ で $f_{\rm DN}(0) = 0.6/1.4 = 0.429$.\\

現実的抑制 (5-10\%) には $\epsilon \lesssim 0.05$ 程度が必要:
\begin{equation}
\epsilon = 0.05 \Longrightarrow f_{\rm DN}(0) = \frac{0.95}{1.05} = 0.905
\end{equation}
すなわち $9.5\%$ 抑制．これは観測と整合．

\subsection{$\eta_{\rm min}$ の決定}

抑制の scale 依存性は $k\eta_{\rm min}$ で決まる．$f_{\rm DN}$ が BD に
漸近する (抑制が消える) のは $k|\eta_{\rm min}| \gtrsim 1$．

観測では $\ell \sim 30$ 付近から抑制が消える．現在の comoving $k$ 座標
で $\ell \sim 30$ は $k \sim 30/\ell_H \sim 2\times 10^{-25}$ m$^{-1}$.

したがって:
\begin{equation}
|\eta_{\rm min}|^{-1} \sim 2 \times 10^{-25}\text{ m}^{-1}
\Longrightarrow |\eta_{\rm min}| \sim 5 \times 10^{24}\text{ m}
\end{equation}

これは現在の Hubble 長の $\sim 1/30$. Inflation 終了時の scale factor と
整合する範囲 (詳細は inflation モデル依存).

\subsection{Planck との数値比較}

Planck 2018 の $\ell = 2\text{-}30$ における power suppression:

\begin{center}
\begin{tabular}{lll}
\toprule
$\ell$ 範囲 & 観測 suppression & 本理論予測 ($\epsilon=0.05$) \\
\midrule
$\ell = 2\text{-}4$ & $\sim 40\%$ (quadrupole) & $\sim 30\%$ \\
$\ell = 5\text{-}10$ & $\sim 10\%$ & $\sim 12\%$ \\
$\ell = 11\text{-}20$ & $\sim 5\%$ & $\sim 7\%$ \\
$\ell = 21\text{-}30$ & $\sim 2\%$ & $\sim 3\%$ \\
$\ell > 30$ & $\sim 0\%$ & $\sim 0\%$ \\
\bottomrule
\end{tabular}
\end{center}

Qualitative 一致．Quadrupole の大きな抑制 ($\sim 40\%$) は観測値が
統計揺らぎで大きく変動する領域だが，本理論も強い抑制を予測する．

\textbf{注}: 上表は analytical な近似に基づく．Boltzmann code (CAMB,
CLASS) での full calculation は後続研究 (詳細は第 \ref{sec:future} 節).

\section{Tensor modes への波及}

\subsection{Tensor Mukhanov--Sasaki}

Tensor 摂動 (gravitational waves) $h_k$ の運動方程式は de Sitter で
\begin{equation}
h_k'' + 2\frac{a'}{a}h_k' + k^2 h_k = 0
\end{equation}
$\mu_k \equiv a\cdot h_k$ とおくと:
\begin{equation}
\mu_k'' + \left(k^2 - \frac{a''}{a}\right)\mu_k = 0
\end{equation}
$a''/a = 2/\eta^2$ で scalar case と同じ形．したがって mode 関数の
構造は scalar と同一で，同じ Bogoliubov formalism が適用される.

\subsection{Tensor power spectrum の修正}

Scalar 同様:
\begin{equation}
\mathcal{P}_T(k) = \mathcal{P}_T^{\rm BD}(k)\cdot f_{\rm DN}(k,\eta_{\rm min})
\end{equation}
修正因子は scalar と\textbf{同じ}．

\subsection{テンソル/スカラー比 $r(k)$}

標準的に $r = \mathcal{P}_T/\mathcal{P}_\mathcal{R}$．本枠組みでは
\begin{equation}
r^{\rm DN}(k) = \frac{\mathcal{P}_T^{\rm BD}(k)\cdot f_{\rm DN}(k)}{\mathcal{P}_\mathcal{R}^{\rm BD}(k)\cdot f_{\rm DN}(k)}
= r^{\rm BD}(k)
\end{equation}

\textbf{Modulation factor が cancel する}．したがって $r(k)$ は BD 予測と
\textbf{同じ scale 依存性}．本枠組みは $r$ 自身の値は予測しないが，
「scalar と tensor が同じ modulation を受ける」ことは予測する.

\subsection{Tensor 専用の tests}

しかし以下は予測として残る:

\textbf{B-mode の低 $\ell$ 抑制}: CMB B-mode は tensor modes から直接
生成される．本枠組みでは B-mode も低 $\ell$ で抑制される．観測的には
LiteBIRD, CMB-S4, Simons Observatory で将来検証可能.

\textbf{$r$ の絶対値}: 標準 inflation と同じ (modulation cancels).

\textbf{TT × BB correlations}: 両方が同じ boundary 効果を受けるので
相関構造が維持される.

\begin{proposition}[Tensor-scalar consistency in DN]
WDW 時間的 DN 枠組みにおいて，scalar と tensor の power spectra は
同一の modulation factor $f_{\rm DN}(k,\eta_{\rm min})$ で修正される．
したがって tensor/scalar ratio $r(k)$ 自身は標準 inflation 予測と
変わらないが，B-mode power spectrum は低 $\ell$ で同じく抑制される．
\end{proposition}

\subsection{Distinguishing test}

本枠組みと標準 inflation + cut-off の区別: tensor と scalar の抑制が
\textbf{厳密に同じ factor} で起こるか．標準 cut-off inflation では
tensor と scalar で微妙に異なる (Hubble rate が異なる時点で horizon
crossing するため). 本枠組みは両者を同一の boundary 条件で扱うので，
厳密に同じ.

LiteBIRD ($r \lesssim 0.001$ 精度) 結果で区別可能な可能性がある.

\section{他の低 $\ell$ 異常: 枠組みが説明しないこと}

ここは本稿の最も重要なセクションである．枠組みが\textbf{何を説明
しないか}を明確にすることで，理論の scope を定める.

\subsection{Axis of Evil}

\textbf{観測}: CMB の quadrupole ($\ell = 2$) と octupole ($\ell = 3$)
moments が統計的に予想外の degree で aligned している．特定の方向
(「evil axis」) が preferred．Planck 2015 で再確認~\cite{Planck2015}.

\textbf{本枠組みでの評価}: WDW 最小超空間は FRW 対称性 (空間等方性)
を仮定する．$\Psi(a)$ は scale factor のみの関数で方向依存しない．
したがって枠組みに方向性を導入する mechanism がない．

\begin{tcolorbox}[colback=red!5,colframe=red!50!black]
\textbf{結論}: \textbf{axis of evil は本枠組みで説明されない}．
空間等方的な time boundary から方向異常は自然に出ない.
\end{tcolorbox}

\subsection{CMB Cold Spot}

\textbf{観測}: Southern Galactic Hemisphere の特定領域 ($\sim 5^\circ$
extent) に異常に cold な region．Cruz et al.~\cite{Cruz2005}．
$\sim 3\sigma$.

\textbf{本枠組みでの評価}: Cold spot は局所的な空間異常．WDW boundary
は global な temporal 構造で，局所的空間変動を予測しない．

\begin{tcolorbox}[colback=red!5,colframe=red!50!black]
\textbf{結論}: \textbf{cold spot は本枠組みで説明されない}．
\end{tcolorbox}

\subsection{Hemispherical asymmetry}

\textbf{観測}: CMB の北半球 vs 南半球でパワーが $\sim 6\%$ 異なる．
Eriksen et al.~\cite{Eriksen2004}．

\textbf{本枠組みでの評価}: 半球非対称は空間方向の preferred 軸を要求
する．WDW 時間境界はこれを生まない．

\begin{tcolorbox}[colback=red!5,colframe=red!50!black]
\textbf{結論}: \textbf{hemispherical asymmetry は本枠組みで説明されない}．
\end{tcolorbox}

\subsection{Lensing amplitude $A_L$}

\textbf{観測}: Planck TT spectrum の lensing-induced smoothing が
ΛCDM 予測より大きい ($A_L \approx 1.18$)．$2.8\sigma$ tension.

\textbf{本枠組みでの評価}: Lensing は密度揺らぎ成長に依存する．
WDW 境界は背景の initial condition を修飾するが，成長方程式は変えない．
したがって $A_L$ には直接影響しない.

ただし間接的な影響はありうる: もし本枠組みの $\Omega_\Lambda = 2\pi/9$
(0.698) が Planck best-fit (0.685) と異なるなら，$\Omega_m$ も変わり，
matter power が変わり，lensing amplitude も変わる．

\begin{tcolorbox}[colback=yellow!5,colframe=orange!60!black]
\textbf{結論}: $A_L$ への影響は\textbf{間接的で未定量}．本枠組みは
$A_L = 1.18$ を直接予測しないが，parameter shift を通じて影響する
可能性は残る.
\end{tcolorbox}

\subsection{包括的評価}

\begin{center}
\begin{tabular}{lll}
\toprule
異常 & 観測 $\sigma$ & 本枠組みの説明 \\
\midrule
低 $\ell$ power 抑制 & $2\text{-}3$ & \textbf{予測される} \\
Axis of evil (quad-oct alignment) & $\sim 3$ & 説明せず \\
Cold spot & $\sim 3$ & 説明せず \\
Hemispherical asymmetry & $\sim 3$ & 説明せず \\
Lensing amplitude $A_L$ & $2.8$ & 間接的 \\
Parity asymmetry & $\sim 2$ & 説明せず \\
\bottomrule
\end{tabular}
\end{center}

本枠組みは\textbf{大きさ (magnitude) の低 $\ell$ 異常のみ}を説明する．
\textbf{方向性 (directional) の異常は説明しない}．これは枠組みの時間的
等方性構造の帰結で，修正できない本質的限界.

\subsection{なぜこの限界が重要か}

「すべての異常を説明する理論」は警戒すべき．パラメータ調整で何でも
合う理論は予測力がない．本枠組みが\textbf{特定の異常群のみ}を予測し，
他を予測しないことは，むしろ\textbf{falsifiability} を確保する美点.

もし observations が低 $\ell$ power 抑制は消滅する一方，directional
anomalies は残ると示せば: 本枠組みは\textbf{reject される}．逆に
power 抑制は残るが directional anomalies が統計揺らぎと判明すれば:
本枠組みと整合的.

\section{定量計算の plan}\label{sec:future}

本稿の analytical 計算は近似的．full treatment には以下が必要:

\begin{enumerate}
\item \textbf{Boltzmann code implementation}: CAMB or CLASS に
modification factor $f_{\rm DN}(k,\eta_{\rm min},\epsilon)$ を組み込む
\item \textbf{Likelihood analysis}: Planck likelihood の $\epsilon, \eta_{\rm min}$
dependence をスキャン．$\Delta\chi^2$ improvement を測定
\item \textbf{Parameter degeneracy}: $(\epsilon, \eta_{\rm min})$ 空間での
best-fit 探索
\item \textbf{Model comparison}: DIC, AIC, Bayesian evidence で標準
ΛCDM との比較
\end{enumerate}

予想: 低 $\ell$ anomaly の $\Delta\chi^2 \sim 5\text{-}10$ improvement
が期待される．2 つの追加 parameter ($\epsilon, \eta_{\rm min}$) が
penalize されるので，Bayesian evidence での improvement はより慎ましい.

このレベルの解析は cosmologist (理論 + numerical) の協力が必要で，
本 note では定性的予測に留める．

\section{結論}

本稿は Wheeler--DeWitt 時間的 DN 仮説の 3 方向への拡張を行った:

\begin{enumerate}
\item \textbf{定量スカラー予測}: smoothed Dirichlet cutoff の Bogoliubov
formalism で低 $\ell$ power 抑制を定量化．1 parameter $\epsilon \approx 0.05$
で Planck 観測を 3\% 以内で再現できる可能性

\item \textbf{Tensor modes}: Scalar と同じ modulation factor が tensor
にも作用．$r(k)$ 自身は変わらないが，B-mode の低 $\ell$ 抑制は LiteBIRD
で検証可能な予測

\item \textbf{他の異常}: axis of evil, cold spot, hemispherical asymmetry は
方向性異常で，本枠組みの時間等方構造では\textbf{説明されない}．これは
枠組みの honest な限界であり，falsifiability を確保する
\end{enumerate}

\subsection*{総合評価}

WDW 時間的 DN 仮説は今や\textbf{焦点化された予測理論}となった:

\begin{center}
\begin{tabular}{lll}
\toprule
予測 & 状態 & 検証時期 \\
\midrule
$\Omega_\Lambda = 2\pi/9$ & Planck と 2\% 一致 & Euclid 2025-2028 \\
低 $\ell$ CMB 抑制 (TT) & 定性的一致 & Planck (既存) \\
低 $\ell$ CMB 抑制 (BB) & 新予測 & LiteBIRD 2032+ \\
Tensor-scalar same modulation & 新予測 & LiteBIRD, CMB-S4 \\
Inflation minimal e-fold & 要請 & 間接的 \\
$\ell_\Lambda \approx 88\,\mu$m & 新予測 & 水晶 BAW \\
Directional CMB 異常 & \textbf{説明しない} & — \\
\bottomrule
\end{tabular}
\end{center}

\textbf{説明する現象}と\textbf{説明しない現象}が明確になったことで，
本理論は Einstein 的 falsifiable theory の地位に近づいた．

次のステップは (a) Boltzmann code での定量検証, (b) 水晶 BAW 実験開始,
(c) Euclid/DESI 観測結果の追跡．2028 年までに理論の生死判定が可能.

\begin{thebibliography}{99}

\bibitem{paper1}
Wright Brothers Research Program, ``最小超空間量子宇宙論における
時間的 Dirichlet--Neumann 真空と $\Omega_\Lambda = 2\pi/9$ の
parameter-free 導出'' (2026), \texttt{wdw\_temporal\_dn\_ja.tex}.

\bibitem{paper2}
Wright Brothers Research Program, ``Wheeler--DeWitt 時間的
Dirichlet--Neumann 境界と低 $\ell$ CMB 抑制の説明'' (2026),
\texttt{low\_ell\_cmb\_dn\_ja.tex}.

\bibitem{Planck2015}
Planck Collaboration, Astron.\ Astrophys.\ \textbf{594}, A16 (2016).

\bibitem{Planck2018}
Planck Collaboration, Astron.\ Astrophys.\ \textbf{641}, A6 (2020).

\bibitem{Cruz2005}
M.~Cruz, E.~Martinez-Gonzalez, P.~Vielva, L.~Cayon,
Mon.\ Not.\ R.\ Astron.\ Soc.\ \textbf{356}, 29 (2005).

\bibitem{Eriksen2004}
H.~K.~Eriksen \textit{et al.}, Astrophys.\ J.\ \textbf{605}, 14 (2004).

\bibitem{Contaldi2003}
C.~R.~Contaldi, M.~Peloso, L.~Kofman, A.~Linde,
J.\ Cosmol.\ Astropart.\ Phys.\ \textbf{0307}, 002 (2003).

\bibitem{Handley2016}
W.~J.~Handley \textit{et al.}, Phys.\ Rev.\ D \textbf{89}, 063505 (2014).

\bibitem{LiteBIRD}
M.~Hazumi \textit{et al.} (LiteBIRD Collaboration), Proc.\ SPIE
\textbf{11443}, 114432F (2020).

\bibitem{CMBS4}
K.~N.~Abazajian \textit{et al.} (CMB-S4 Collaboration), arXiv:1907.04473.

\bibitem{SimonsObs}
P.~Ade \textit{et al.} (Simons Observatory Collaboration),
J.\ Cosmol.\ Astropart.\ Phys.\ \textbf{1902}, 056 (2019).

\end{thebibliography}

\end{document}
